% Manuscript-ready section. Requires amsmath, amssymb, amsthm.
% Bibliographic keys: ghp:ADH, ghp:DLMF, ghp:ExactResonant.
\section{Exact parity of generalized harmonic numerators}
\label{ghp:sec}

Question Q11 in the pinned \emph{Exact Resonant Identities} report
\cite{ghp:ExactResonant} asks which polynomials in generalized harmonic
numbers have no principal part at any nonpositive even spectral integer
when the outer shift is $1/2$. The raw-power argument proves a sufficient
condition for powers of $H_n$; it does not itself classify arbitrary
polynomials. This section gives a necessary and sufficient algebraic
criterion. The necessity uses an established theorem on formal
difference equations. It does not assume algebraic independence of
Euler's constant, zeta values, or numerical periods.

\subsection{The centered involution and the classification theorem}

Put $H_0^{(r)}=0$, $H_n^{(r)}=\sum_{k=1}^n k^{-r}$, and $H_n=H_n^{(1)}$.
For a fixed $R\geq1$ and a polynomial $P\in\mathbb C[X_1,\ldots,X_R]$,
let
\begin{equation}\label{ghp:def:series}
 \mathcal D_P(s,a)=\sum_{n=0}^{\infty}
  \frac{P(H_n,H_n^{(2)},\ldots,H_n^{(R)})}{(n+a)^s},
 \qquad a>0,\quad \Re s>1.
\end{equation}
These series converge normally in that half-plane; a polynomial
numerator grows at most as a fixed power of $\log(n+2)$.
All later values refer to the unique meromorphic continuation.

Define an involution on the polynomial ring by
\begin{equation}\label{ghp:def:tau}
 \tau X_1=X_1,\qquad
 \tau X_r=\begin{cases}
       2\zeta(r)-X_r,&r\geq2\text{ even},\\
       X_r,&r\geq3\text{ odd}.
       \end{cases}
\end{equation}
The coefficients in $\mathbb C$ are fixed. Thus the transformation
changes the sign of every even-order tail $\zeta(2j)-X_{2j}$
\emph{simultaneously}.

\Needspace{20\baselineskip}
\begin{theorem}[Exact centered cancellation criterion]\label{ghp:thm:classification}
For $P\in\mathbb C[X_1,\ldots,X_R]$, the following conditions are
 equivalent:
\begin{enumerate}
\item $\mathcal D_P(s,1/2)$ is holomorphic at every
      $s=0,-2,-4,\ldots$.
\item $\tau P=P$ as a polynomial with complex coefficients.
\item On writing $V_j=\zeta(2j)-X_{2j}$, every monomial of $P$ has
      even total degree in the variables $V_1,V_2,\ldots$.
\end{enumerate}
Equivalently, the complete algebra of such numerators is
\begin{equation}\label{ghp:eq:invariant-ring}
 \mathbb C\bigl[X_1,X_3,X_5,\ldots,
  (\zeta(2i)-X_{2i})(\zeta(2j)-X_{2j}):1\leq i\leq j\leq\lfloor R/2\rfloor\bigr].
\end{equation}
Only variables whose index is at most $R$ are retained.
\end{theorem}

The paired generators in \eqref{ghp:eq:invariant-ring} are a generating
set, not an assertion that this algebra is freely generated. For
example, $(V_iV_j)^2=V_i^2V_j^2$ is an algebraic relation among them.

Examples of admissible numerators are
\[
 H_n^pH_n^{(3)},\qquad
 (\zeta(2)-H_n^{(2)})^2,\qquad
 H_n^p(\zeta(2)-H_n^{(2)})(\zeta(4)-H_n^{(4)}),
 \qquad p\geq0.
\]
The raw polynomial $H_n^{(2)}$ does not qualify. In fact its principal
part at $s=0$ is $-1/s$, since
$H_n^{(2)}=\zeta(2)-(n+1/2)^{-1}+O(n^{-3})$.
Multiplying two even-order tails changes this outcome by making their
full centered expansion even.

\subsection{A finite formula for every principal part}

Let $L$ be an independent formal symbol, $z$ a formal variable, and
\begin{equation}\label{ghp:def:A}
 A(z,a)=-\sum_{k\geq1}\frac{B_k(a)}kz^k,
 \qquad T_1(z,L;a)=L+\gamma+A(z,a).
\end{equation}
For $r\geq2$ define
\begin{equation}\label{ghp:def:Tr}
 T_r(z;a)=\zeta(r)-\frac{z^{r-1}}{r-1}
 -\frac1{r-1}\sum_{k\geq1}
   \binom{k+r-2}{r-2}B_k(a)z^{k+r-1}.
\end{equation}
These are formal asymptotic series; convergence in $z$ is not needed.
Set
\begin{equation}\label{ghp:def:q}
 P(T_1,T_2,\ldots,T_R)=\sum_{j\geq0}q_{P,j}(L;a)z^j,
 \qquad q_{P,j}(L;a)=\sum_{\ell\geq0}q_{P,j,\ell}(a)L^\ell.
\end{equation}
Every coefficient requires only finitely many Bernoulli polynomials and
finite polynomial operations.

\begin{proposition}[Generalized-harmonic principal coefficients]
\label{ghp:prop:principal}
The function \eqref{ghp:def:series} continues meromorphically to the
entire $s$-plane. Its possible poles are $1,0,-1,-2,\ldots$. At
$s=1-j$ its complete principal part is
\begin{equation}\label{ghp:eq:principal}
 \sum_{\ell\geq0}\frac{\ell!\,q_{P,j,\ell}(a)}{(s+j-1)^{\ell+1}}.
\end{equation}
The sum is finite. In particular, this point is regular if and only if
$q_{P,j}(L;a)$ is the zero polynomial in $L$.
\end{proposition}

\begin{proof}
Writing $x=n+a$, the ordinary shifted digamma expansion and its fixed
argument derivatives give the expansions $T_r$ above, with $z=x^{-1}$
and $L=\log x$; see the classical Gamma and polygamma formulas in
\cite[\S\S5.11, 5.15]{DLMF}. For $r\geq2$ one can also obtain
\eqref{ghp:def:Tr} directly by differentiating
\[
 H_n=\psi(n+1)+\gamma,
 \qquad
 \zeta(r)-H_n^{(r)}=
       \frac{(-1)^r}{(r-1)!}\psi^{(r-1)}(n+1).
\]
At every fixed truncation order $J$, the polynomial substitution has
remainder
\[
 O\bigl(x^{-J-1}(1+\log^d x)\bigr)
\]
for some fixed $d$.
Subtract those finitely many terms from the Dirichlet series. The
remaining series converges normally in $\Re s>-J$ on compact subsets.
The subtracted terms are
\[
 \sum_{j=0}^J\sum_\ell
   q_{P,j,\ell}(a)(-1)^\ell\partial_s^\ell\zeta(s+j,a).
\]
Their only principal contribution at $s=1-j$ is
\eqref{ghp:eq:principal}. Taking arbitrarily large $J$ proves the
continuation. Distinct powers $(s+j-1)^{-\ell-1}$ cannot cancel one
another, proving the last assertion.
\end{proof}

At $a=1/2$, all odd Bernoulli polynomials vanish. Consequently
\begin{equation}\label{ghp:eq:formal-parity}
 T_1(-z,L;1/2)=T_1(z,L;1/2),\qquad
 T_r(-z;1/2)=\begin{cases}
  2\zeta(r)-T_r(z;1/2),&r\text{ even},\\
  T_r(z;1/2),&r\text{ odd}.
 \end{cases}
\end{equation}
This proves the sufficiency of $P=\tau P$. The necessity requires
showing that no nonzero polynomial is annihilated by the formal
substitution. A finite list of verified asymptotic coefficients could
not establish that fact.

\subsection{Why the formal substitution is injective}

We use the following established result in exactly its formal Laurent
setting. Adamczewski, Dreyfus, and Hardouin
\cite[Theorem~1.2, case $\mathbf S_\infty$]{ghp:ADH} prove that a formal
Laurent series $f\in\mathbb C((x^{-1}))$ satisfying a homogeneous linear
difference equation over $\mathbb C(x)$ for the shift $x\mapsto x+1$
is either rational or differentially transcendental over
$\mathbb C(x)$. The latter means that every finite tuple
$f,f',\ldots,f^{(q)}$ is algebraically independent over $\mathbb C(x)$.

\begin{lemma}[Independence of the centered formal harmonic tails]
\label{ghp:lem:injective}
The substitution
\[
 \mathbb C[X_1,\ldots,X_R]\longrightarrow\mathbb C[L][[z]],
 \qquad X_r\longmapsto T_r(z,L;1/2),
\]
is injective.
\end{lemma}

\begin{proof}
Let
\begin{equation}\label{ghp:eq:formal-trigamma}
 f(x)=\frac1x+\sum_{k\geq1}B_{2k}(1/2)x^{-2k-1}
        \in\mathbb C((x^{-1})).
\end{equation}
This is the formal expansion of $\psi'(x+1/2)$. The trigamma shift
identity, or the Bernoulli generating series, gives
\begin{equation}\label{ghp:eq:shift}
 f(x+1)-f(x)=-\frac1{(x+1/2)^2}.
\end{equation}
In particular it satisfies the homogeneous second-order equation
\begin{equation}\label{ghp:eq:homogeneous}
 f(x+2)-(1+c(x))f(x+1)+c(x)f(x)=0,
 \qquad c(x)=\frac{(x+1/2)^2}{(x+3/2)^2}.
\end{equation}

There is no rational solution of \eqref{ghp:eq:shift}. Indeed, in any
nonempty finite set of poles of a rational function belonging to one
coset modulo $\mathbb Z$, choose the first and last pole in the integer
ordering. Its forward difference has an uncancelled pole immediately
before the first and another at the last. These poles are distinct.
Thus a rational forward difference cannot have precisely one pole in
such a coset. The right side of \eqref{ghp:eq:shift} has exactly one
pole. A rational function with no finite poles is a polynomial, whose
forward difference is polynomial, and is also impossible.

The cited difference-equation theorem applied to
\eqref{ghp:eq:homogeneous} now shows that
$f,f',\ldots,f^{(R-2)}$ are algebraically independent over
$\mathbb C(x)$ when $R\geq2$. For $r\geq2$, the substitution for $X_r$
is the affine expression
\[
 T_r=\zeta(r)-\frac{(-1)^r}{(r-1)!}f^{(r-2)}.
\]
These tails are therefore algebraically independent, even over
$\mathbb C(x)$. Finally $T_1=L+\gamma+A(z,1/2)$ has coefficient one
in the independent formal symbol $L$, and the other $T_r$ contain no
$L$. If a nonzero polynomial in $X_1$ has degree $d$ and nonzero leading
coefficient $a_d(X_2,\ldots,X_R)$, the coefficient of $L^d$ after
substitution is the nonzero substituted value of $a_d$. This proves
injectivity, including the case $R=1$ directly.
\end{proof}

\begin{proof}[Proof of Theorem~\ref{ghp:thm:classification}]
By Proposition~\ref{ghp:prop:principal}, regularity at every $s=-2m$
is equivalent to the vanishing of every odd coefficient $q_{P,2m+1}$.
Equivalently, the complete formal substitution is even in $z$ while
$L$ is held fixed. By \eqref{ghp:eq:formal-parity}, this says that the
substitution annihilates $P-\tau P$. Lemma~\ref{ghp:lem:injective}
therefore gives $P-\tau P=0$.

Writing the even-index variables as $\zeta(2j)-V_j$ makes $\tau$ the
simultaneous sign change of all the $V_j$. The invariant polynomials
are exactly those with even total degree in these variables. Every
such monomial can be grouped into pairs $V_iV_j$, proving the stated
finite generating description.
\end{proof}

\begin{corollary}[Complementary pole lattices]\label{ghp:cor:split}
Write $P=P_++P_-$, where $P_\pm=(P\pm\tau P)/2$. Then
$\mathcal D_{P_+}(s,1/2)$ has possible poles only at
$1,-1,-3,\ldots$, while $\mathcal D_{P_-}(s,1/2)$ has possible poles
only at $0,-2,-4,\ldots$. Moreover, if $P\ne0$,
$\mathcal D_P(s,1/2)$ is not entire.
\end{corollary}
\begin{proof}
The corresponding formal expansions have only even and only odd
powers of $z$, respectively. This gives the two pole lists by
\eqref{ghp:eq:principal}. If the whole Dirichlet function were entire,
all its coefficient polynomials $q_{P,j}$ would vanish. Injectivity
would then force $P=0$.
\end{proof}

This is a statement about identities of functions and formal series,
not about algebraic independence of their special numerical values.
The criterion concerns cancellation at \emph{every} point of the
specified parity. At any finite list of spectral points, additional
accidental cancellations can occur and are decided by the finite
Bernoulli conditions \eqref{ghp:eq:principal}.

\section{All centered trigamma-square moments and a polylogarithmic generator}
\label{ghp:square:sec}

The invariant numerator $(\zeta(2)-H_n^{(2)})^2$ provides an explicit
family beyond raw powers of $H_n$. Define
\begin{equation}\label{ghp:def:square}
 A_n=\zeta(2)-H_n^{(2)}=\psi'(n+1),\qquad
 \mathcal T(s)=\sum_{n=0}^{\infty}\frac{A_n^2}{(n+1/2)^s}.
\end{equation}
The defining series already converges for $\Re s>-1$. By the preceding
theorem it is regular at every nonpositive even integer. All its values
there reduce to a rational number plus a rational multiple of
$\zeta(3)$.

\begin{theorem}[Every nonpositive even trigamma-square value]
\label{ghp:thm:square}
For every integer $r\geq0$,
\begin{align}\label{ghp:eq:all-square}
 \mathcal T(-2r)={}&3B_{2r}(1/2)\zeta(3)\notag\\
 &-\frac12\sum_{j=1}^{r}\binom{2r}{2j}
   \frac{2j-1}{2j+1}B_{2r-2j}(1/2)B_{2j-2}.
\end{align}
The Bernoulli number convention is $B_1=-1/2$, and the sum is empty
when $r=0$.
\end{theorem}

For example,
\begin{align}\label{ghp:eq:square-table}
 \mathcal T(0)&=3\zeta(3),\notag\\
 \mathcal T(-2)&=-\frac14\zeta(3)-\frac16,\notag\\
 \mathcal T(-4)&=\frac7{80}\zeta(3)+\frac1{30},\notag\\
 \mathcal T(-6)&=-\frac{31}{448}\zeta(3)+\frac1{672},\notag\\
 \mathcal T(-8)&=\frac{127}{1280}\zeta(3)-\frac{17}{540}.
\end{align}
The first value is the classical Euler sum. The theorem supplies a
single finite formula for all subsequent centered moments.

\subsection{Equivalent identities of absolutely convergent ordinary sums}

Let
\begin{equation}\label{ghp:def:ck}
 c_k=\sum_{i=0}^{k-1}B_{2i}(1/2)B_{2k-2-2i}(1/2),\qquad k\geq1.
\end{equation}
Thus $c_1=1$, $c_2=-1/6$, $c_3=47/720$, and
$c_4=-257/5040$. Squaring \eqref{ghp:eq:formal-trigamma} gives
$A_n^2\sim\sum_{k\geq1}c_k(n+1/2)^{-2k}$.
For every $r\geq0$ the exact ordinary convergent identity is
\begin{equation}\label{ghp:eq:convergent-square}
 \sum_{n=0}^{\infty}\left\{(n+1/2)^{2r}\psi'(n+1)^2
       -\sum_{k=1}^r c_k(n+1/2)^{2r-2k}\right\}
       =\mathcal T(-2r),
\end{equation}
where the right side is explicitly \eqref{ghp:eq:all-square}. The
summands after subtraction are $O(n^{-2})$.

For instance,
\[
 \sum_{n=0}^{\infty}\left((n+1/2)^2\psi'(n+1)^2-1\right)
       =-\frac14\zeta(3)-\frac16,
\]
and
\[
 \sum_{n=0}^{\infty}\left((n+1/2)^4\psi'(n+1)^2
                         -(n+1/2)^2+\frac16\right)
       =\frac7{80}\zeta(3)+\frac1{30}.
\]
These statements need no prescription for summing a divergent series.
To prove \eqref{ghp:eq:convergent-square}, subtract the displayed
asymptotic terms in the half-plane of convergence and continue. Their
Dirichlet sums at $s=-2r$ are
$c_k\zeta(2k-2r,1/2)=0$, since
$\zeta(-2m,1/2)=0$ for $m\geq0$.

\subsection{Proof of the Bernoulli formula by finite summation}

For a sequence admitting an expansion in powers of $N$ and $\log N$,
write $\operatorname{CT}_N$ for its coefficient of $N^0(\log N)^0$.
Only finite exact identities followed by known asymptotic expansions
will be used here. Put
\begin{equation}\label{ghp:def:S}
 S_r(k)=\sum_{n=0}^{k-1}(n+1/2)^{2r}
  =\frac{B_{2r+1}(k+1/2)}{2r+1}
  =\sum_{j=0}^r b_{r,j}k^{2j+1},
\end{equation}
where
\begin{equation}\label{ghp:def:b}
 b_{r,j}=\frac{\binom{2r}{2j}}{2j+1}B_{2r-2j}(1/2).
\end{equation}
The Bernoulli polynomial $B_{2r+1}(1/2)$ is zero, so this is an odd
polynomial in $k$ and has zero constant term.

\begin{proof}[Proof of Theorem~\ref{ghp:thm:square}]
First, \eqref{ghp:eq:convergent-square} and \eqref{ghp:def:S} show that
\begin{equation}\label{ghp:eq:ct-reduction}
 \mathcal T(-2r)=\operatorname{CT}_N
       \sum_{n=0}^{N-1}(n+1/2)^{2r}A_n^2.
\end{equation}
Indeed every subtracted even power has partial sum with zero constant
term in $N$, and the convergent remainder tends to the stated value.

The identity $A_{k-1}=A_k+k^{-2}$ gives finite summation by parts:
\begin{align}\label{ghp:eq:partial-sum}
 \sum_{n=0}^{N-1}(n+1/2)^{2r}A_n^2
 ={}&\sum_{k=1}^N S_r(k)
          \left(\frac{2A_k}{k^2}+\frac1{k^4}\right)
       +S_r(N)A_N^2.
\end{align}
The coefficient $b_{r,0}$ of $k$ contributes
\[
 \operatorname{CT}_N
 \left(2\sum_{k=1}^N\frac{A_k}{k}+H_N^{(3)}+NA_N^2\right)
 =3\zeta(3).
\]
Here $\sum_{k\geq1}A_k/k=\zeta(2,1)=\zeta(3)$ is the classical
Euler identity and $NA_N^2\to0$. For completeness, the Euler identity
follows from
\[
 \zeta(2,1)=\frac12\int_0^1\frac{\log^2(1-t)}t\,dt
 =\frac12\int_0^\infty\frac{u^2}{e^u-1}\,du=\zeta(3).
\]

For $j\geq1$, it remains to compute the constant of
\begin{equation}\label{ghp:def:Ej}
 E_j(N)=2\sum_{k=1}^Nk^{2j-1}A_k
        +\sum_{k=1}^Nk^{2j-3}+N^{2j+1}A_N^2.
\end{equation}
Let
\[
 F_j(N)=\sum_{k=1}^Nk^{2j-1}
       =\frac{B_{2j}(N+1)-B_{2j}}{2j}.
\]
Changing the order in a finite sum gives
\begin{equation}\label{ghp:eq:finite-change-order}
 \sum_{k=1}^Nk^{2j-1}A_k
  =A_NF_j(N)+\sum_{\ell=1}^N\frac{F_j(\ell-1)}{\ell^2}.
\end{equation}
For $j\geq2$, the second summand is a sum of polynomial powers of
$\ell$ and has zero constant term as a polynomial in $N$;
$B_{2j-1}=0$ removes a possible $1/\ell$ term. The other power sum in
\eqref{ghp:def:Ej} also has zero constant. For $j=1$ the second sum in
\eqref{ghp:eq:finite-change-order} equals $(N-H_N)/2$; its harmonic
term cancels exactly against $\sum k^{-1}$ in \eqref{ghp:def:Ej}.
In all cases,
\begin{equation}\label{ghp:eq:Ej-CT}
 \operatorname{CT}_NE_j(N)
   =2\operatorname{CT}_N(A_NF_j(N))
          +\operatorname{CT}_N(N^{2j+1}A_N^2).
\end{equation}

Now use the uncentered tail expansion
\[
 A_N=\frac1N-\frac1{2N^2}
                  +\sum_{k\geq1}\frac{B_{2k}}{N^{2k+1}}.
\]
The finite coefficient calculation gives
\begin{equation}\label{ghp:eq:two-constants}
 \operatorname{CT}_N(A_NF_j(N))
      =\frac{3-2j}{4}B_{2j-2},\qquad
 \operatorname{CT}_N(N^{2j+1}A_N^2)=-B_{2j-2}.
\end{equation}
To check the first identity, the coefficient of $N^\ell$ in $F_j$
is $\binom{2j}{\ell}B_{2j-\ell}(1)/(2j)$. Only $\ell=2$ and
$\ell=2j-1$ can pair with nonzero inverse coefficients of $A_N$
when $j\geq2$; their contributions are
$-(2j-1)B_{2j-2}/4$ and $B_{2j-2}/2$. The case $j=1$ is the direct
product with $F_1(N)=(N^2+N)/2$ and gives $1/4$.
For the second identity, an odd inverse exponent in $A_N^2$ arises
only by multiplying $-1/(2N^2)$ by the appropriate odd inverse power;
this gives $-B_{2j-2}$, also at $j=1$.

Combining \eqref{ghp:eq:Ej-CT} and \eqref{ghp:eq:two-constants} yields
\[
 \operatorname{CT}_NE_j(N)=-\frac{2j-1}{2}B_{2j-2}.
\]
Multiplication by $b_{r,j}$, followed by the $j=0$ contribution,
proves \eqref{ghp:eq:all-square}.
\end{proof}

\subsection{A convergent generator involving classical polylogarithms}

The complete sequence has an analytic exponential generating function.
Let
\[
 \mathcal G(t)=\sum_{r=0}^{\infty}
               \mathcal T(-2r)\frac{t^{2r}}{(2r)!}.
\]

\begin{theorem}[Polylogarithmic moment generator]\label{ghp:thm:generator}
The series $\mathcal G$ converges for $|t|<2\pi$ and equals
\begin{equation}\label{ghp:eq:integral-generator}
 \mathcal G(t)=\frac{t}{2\sinh(t/2)}
   \left[3\zeta(3)-\frac1{4t}
       \int_0^t(t-u)u^2\coth(u/2)\,du\right].
\end{equation}
The singularities at zero in this expression are removable. For
real $0<t<2\pi$, an equivalent formula in classical polylogarithms is
\begin{align}\label{ghp:eq:polylog-generator}
 \mathcal G(t)=\frac{t}{2\sinh(t/2)}\bigg[
 2\zeta(3)-\frac{t^3}{48}
 -\frac t2\operatorname{Li}_2(e^{-t})
 -2\operatorname{Li}_3(e^{-t})
 +\frac{3\bigl(\zeta(4)-\operatorname{Li}_4(e^{-t})\bigr)}t
 \bigg].
\end{align}
Its continuation specified by \eqref{ghp:eq:integral-generator} is
holomorphic and even throughout $|t|<2\pi$.
\end{theorem}

\begin{proof}
The Bernoulli generating functions give
\[
 \frac{t}{2\sinh(t/2)}
     =\sum_{r\geq0}B_{2r}(1/2)\frac{t^{2r}}{(2r)!},
 \qquad
 \frac u2\coth(u/2)
     =\sum_{k\geq0}B_{2k}\frac{u^{2k}}{(2k)!}.
\]
Both converge in the disk of radius $2\pi$. Integrating the second
series twice gives
\[
 \frac1{2t}\int_0^t(t-u)u^2\coth(u/2)\,du
  =\sum_{j\geq1}\frac{2j-1}{2j+1}B_{2j-2}
                  \frac{t^{2j}}{(2j)!}.
\]
Thus \eqref{ghp:eq:integral-generator} is exactly the binomial
convolution in \eqref{ghp:eq:all-square}. This proves convergence,
analyticity and the integral formula. In particular all apparent odd
or logarithmic terms in any branch representation must cancel.

For positive $t$, expand $(e^u-1)^{-1}$ in its convergent positive
exponential series. Direct integration, justified for the integrable
factor $(t-u)u^2$, gives
\begin{align*}
 \int_0^t\frac{(t-u)u^2}{e^u-1}\,du
  ={}&2t\zeta(3)-6\zeta(4)
      +t^2\operatorname{Li}_2(e^{-t})\\
     &+4t\operatorname{Li}_3(e^{-t})
      +6\operatorname{Li}_4(e^{-t}).
\end{align*}
Finally $\coth(u/2)=1+2/(e^u-1)$ and
$\int_0^t(t-u)u^2\,du=t^4/12$ give
\eqref{ghp:eq:polylog-generator}.
\end{proof}

\subsection{Verification, audit scope, and further exact questions}

The accompanying \texttt{verify\_harmonic\_parity.py} implements the
formal involution, verifies the centered trigamma difference equation,
and checks the two finite Bernoulli constant calculations independently
for eight values of $j$. It evaluates seven spectral values directly
from finite trigamma sums with independently truncated asymptotic tails
at two cutoffs, using 110 decimal digits. The largest discrepancy
against \eqref{ghp:eq:all-square} is less than $1.6\times10^{-85}$.
At three positive parameters it compares the polylogarithmic generator
both with the Bernoulli moment series and with ordinary quadrature of
\eqref{ghp:eq:integral-generator}. These are floating-point diagnostics,
not interval enclosures. The proofs do not depend on their precision.

The selected harmonic statements in the pinned report were consistent
with these checks. No new false formula was identified in that selected
source. Its Q11 warning was correct: raw-power parity alone does not
apply to an arbitrary higher-harmonic numerator. The result here
resolves that specific question by supplying the exact invariant ring;
it does not label the earlier warning an error.

The following are concrete extensions beyond the proved results.
\begin{enumerate}
\item Evaluate every centered moment of
$(\zeta(2i)-H_n^{(2i)})(\zeta(2j)-H_n^{(2j)})$ by a uniform finite
Bernoulli formula, and determine the precise convergent multiple-zeta
coordinates that remain. The current trigamma-square result is
$i=j=1$.
\item For invariant products of three or more tail pairs, seek a
polylogarithmic generator analogous to \eqref{ghp:eq:polylog-generator}
and determine when the resulting multiple zeta values reduce to
ordinary zeta products.
\item Classify cancellation at a \emph{finite} prescribed set of
spectral integers by the finite coefficient map
\eqref{ghp:eq:principal}. Its kernel can be larger than the invariant
ring and is a tractable exact linear problem in any bounded polynomial
space.
\item Derive exact spectral-derivative identities for $\mathcal T$
at its regular even points. Such derivatives retain a global remainder
and do not follow by formally differentiating the discrete index $r$
in \eqref{ghp:eq:all-square}.
\item Determine an elementary proof of the injectivity lemma tailored
to the additive trigamma difference equation. This would remove the
use of the general difference Galois theorem, while leaving the current
unconditional classification unchanged.
\end{enumerate}
